\subsection{逆元}

设 A 是一个幺 Banach 代数，其单位元记作 1。回顾 (TG, IX, 命题 14, p.~40)，A 的可逆元所成的群 G 是 A 的一个开子集，A 的拓扑在 G 上诱导的拓扑与群结构相容，且拓扑群 G 是完备的。

命题 2. --- 设 A 是一个 Banach 代数，x 是 A 的一个元素。级数 \(\sum_{n=1}^{\infty}\lambda^{n}x^{n}\) 作为 \(\lambda\) 的幂级数来考虑，收敛半径为 \(\varrho(x)^{-1}\)。若 A 是幺的且 \(\varrho(x)<1\)，则 1-x 可逆，且其逆为 \(\sum_{n=0}^{\infty}x^{n}\)。

级数 \(\sum_{n=1}^{\infty}\lambda^{n}x^{n}\) 的收敛半径为

\[
(\limsup _ {n \to + \infty} \| x ^ {n} \| ^ {1 / n}) ^ {- 1} = \varrho (x) ^ {- 1}
\]

(参见 VAR, R1, p.~23, 3.1.4)。假定 A 有单位元。若 \(\varrho(x) < 1\)，则级数 \(\sum_{n=0}^{\infty} x^n\) 绝对收敛。由于

\[
(1 - x) \left(\sum_ {n = 0} ^ {k} x ^ {n}\right) = \left(\sum_ {n = 0} ^ {k} x ^ {n}\right) (1 - x) = 1 - x ^ {k + 1}
\]

对每个整数 \(k \geqslant 0\) 成立，元素 1 - x 可逆，且其逆等于 \(\sum_{n=0}^{\infty} x^{n}\)。

推论 1. --- 若 A 是一个幺 Banach 代数，则 A 的可逆元所成的群包含以 1 为中心、1 为半径的开球。

这是直接的，因为 \(\|x\| < 1\) 蕴涵 \(\varrho(x) < 1\)。

推论 2. --- 设 A 是一个 Banach 代数，I 是 A 的一个正则极大左（相应地，右）理想。则 I 是闭的。

设 ( \(\tilde{A}, e\) ) 是由 A 添加单位元而得到的幺 Banach 代数。存在 \(\tilde{A}\) 的一个极大左（相应地，右）理想 J，使得 \(J \cap A = I\) (A, VIII, p.~428, 命题 4)。于是 J 与以 e 为中心、1 为半径的开球不相交 (推论 1)，因而 \(\overline{J} \neq \widetilde{A}\)。由于 J 是极大理想，这就蕴涵 \(\overline{J} = J\)，从而 \(I = J \cap A = \overline{J} \cap A\) 在 A 中是闭的。

推论 3. --- Banach 代数的根是闭的。

事实上，根是诸正则极大左理想之交 (A, VIII, p.~430, 定义 3)。

命题 3. --- 设 A 是一个幺 Banach 代数。

\begin{enumerate}
\def\labelenumi{\alph{enumi})}
\item
  若 \(x \in A\) 有一个左逆（相应地，右逆）y，则每个元素 \(x' \in A\)（满足 \(\|x' - x\| < \|y\|^{-1}\)）都有一个左逆（相应地，右逆）。
\item
  A 的可逆（相应地，左可逆，相应地，右可逆）元素所成的集合在 A 中是开的。
\item
  设 \((x_{n})\) 是 A 的元素的一个序列，它们有左逆（相应地，右逆）\(y_{n}\)，使得 \((x_{n})\) 收敛于一个元素 \(x \in A\)，且序列 \((y_{n})\) 有界。则 \(x\) 是左可逆的（相应地，右可逆的）。
\end{enumerate}

只须处理左逆的情形；右逆的情形通过考虑反代数而得到。

设 \(x, y, x' \in \mathrm{A}\)，使得 \(yx = 1\) 且 \(\|x' - x\| < \|y\|^{-1}\)。我们有

\[
\left\| 1 - y x ^ {\prime} \right\| = \left\| y x - y x ^ {\prime} \right\| \leqslant \left\| y \right\| \cdot \left\| x - x ^ {\prime} \right\| <   1,
\]

因而 \(yx'\) 可逆：存在 \(z \in A\)，使得 \(z(yx') = 1\)。于是元素 \(x'\) 左可逆，其逆为 zy。这就证明了断言 a)，而断言 b) 立即得出。

  设 \((x_{n})\) 是 A 的元素的一个序列，它们有左逆 \(y_{n}\)，使得 \((x_{n})\) 收敛于一个元素 \(x \in A\)，且序列 \((y_{n})\) 有界。若 \(M \geqslant 1\) 是使 \(\|y_{n}\| \leqslant M\)（对一切 \(n \geqslant 1\)）的实数，则对充分大的 n，我们有 \(\|x_{n} - x\| < M^{-1} \leqslant \|y_{n}\|^{-1}\)，因而由 a)，x 有一个左逆。

定义 4. --- 设 A 是一个赋范代数，x 是 A 的一个元素。用 \(\gamma_x\) 与 \(\delta_x\) 表示从 A 到 A 的映射 \(y \mapsto xy\) 与 \(y \mapsto yx\)。我们说 \(x\) 是一个左（相应地，右）拓扑零因子，如果 \(\gamma_x\)（相应地，\(\delta_x\)）不是从 A 到 \(\gamma_x(A)\)（相应地，到 \(\delta_x(A)\)）上的一个同胚。

注. --- 由 TG, IX, p.~36, 推论 2，x 是一个左（相应地，右）拓扑零因子必须且只需存在 A 中的一个序列 \((z_{n})\)，使得 \(\|z_{n}\|=1\)，且使得 \(xz_{n}\)（相应地 \(z_{n}x\)）当 \(n\to+\infty\) 时趋于 0。

左（相应地，右）拓扑零因子是一个左（相应地，右）零因子。假定 A 非零且是幺的。左（相应地，右）拓扑零因子 x 不是左（相应地，右）可逆的。事实上，例如若 yx = 1 且 \(xz_{n}\) 趋于 0，则 \(z_{n} = y(xz_{n})\) 趋于 0，从而不可能对一切 n 有 \(\|z_{n}\| = 1\)。

命题 4. --- 设 A 是一个幺 Banach 代数。设 x 是 A 的一个非左可逆的元素。若存在 A 的左可逆元素的一个序列 \((x_{n})\)，它收敛于 x，则 x 是一个右拓扑零因子。

设 \(y_{n}\) 是 \(x_{n}\) 的一个左逆。由命题 3 (ii)，\(\|y_{n}\|\) 趋于 \(+\infty\)。令 \(z_{n} = \|y_{n}\|^{-1}y_{n}\)。我们有 \(\|z_{n}\| = 1\)，且 \(z_{n}x_{n} = \|y_{n}\|^{-1}\) 趋于 0，因而 \(z_{n}x = z_{n}x_{n} + z_{n}(x - x_{n})\) 趋于 0。于是借助定义 4 后面的注即可得出结论。

命题 5. --- 设 A 是一个幺 Banach 代数，B 是 A 的一个全子代数。则 \(\overline{\mathbf{B}}\) 是 A 的一个全子代数。

事实上，设 \(x\) 是 \(\overline{\mathrm{B}}\) 的一个在 A 中可逆的元素，并设 \((x_n)\) 是 B 的元素的一个序列，它趋于 \(x\)。则对充分大的 \(n\)，\(x_n\) 在 A 中可逆，且 \(x_n^{-1}\) 趋于 \(x^{-1}\)。由于子代数 B 是全的，我们有 \(x_n^{-1} \in \mathrm{B}\)，从而 \(x^{-1} \in \overline{\mathrm{B}}\)。

设 A 是一个幺 Banach 代数，\((y_{i})_{i\in\mathrm{I}}\) 是 A 的元素的一个族。设 B 是由诸元素 \(y_{i}\) 生成的 A 的全子代数。则 \(\overline{B}\) 是包含诸 \(y_{i}\) 的 A 的最小闭全子代数。它称为由诸元素 \(y_{i}\) 生成的闭全子代数。

